\chapter{The Go-Powered Userland}
\label{chap:userland}

Above the C kernel modules and the C \code{au-init} process described in the
previous chapters sits the part of AULinux a human actually touches: the
\emph{userland}. Every interactive command---listing a directory, copying a
file, inspecting running processes, bringing a network interface up---is served
by Go code. There is no GNU coreutils, no Bash, no \code{/bin/sh} inherited from
another distribution. Instead, AULinux ships a single self-contained Go program,
\code{auutils}, that impersonates every core utility, plus an interactive shell,
\code{aush}.

This chapter documents that userland exactly as the source tree implements it.
We start with the design motivation, then walk the multicall dispatcher in
\file{utils/main.go}, tour the implemented coreutils with real code, dissect the
\code{aush} shell line by line, and close with the testing strategy, the static
build pipeline, and an honest accounting of what is and is not finished.

\section{Why Go for the Userland}

The kernel and \code{au-init} are written in C because they run before any
runtime exists and must speak directly to the hardware and the kernel ABI. The
userland has the opposite set of constraints. It is large in surface area
(dozens of commands, lots of string and filesystem manipulation) and it benefits
enormously from a language that closes off entire bug classes. Go was chosen for
three concrete reasons.

\begin{itemize}
  \item \textbf{Static, self-contained binaries.} A Go program compiled with
        \code{CGO\_ENABLED=0} has no dependency on \code{glibc}, \code{ld-linux.so},
        or any shared object. It is a single ELF file that runs in an otherwise
        empty root filesystem. This matches the philosophy established by
        \code{au-init} in Chapter~2, which is itself statically linked. The
        userland never has to worry about a missing \code{.so}.
  \item \textbf{Memory safety without a GC pause problem.} Bounds-checked
        slices, no manual \code{free}, and no dangling pointers eliminate the
        memory-corruption vulnerabilities that have historically plagued
        C reimplementations of coreutils. The garbage collector is a non-issue
        for short-lived command processes, and AULinux tunes it anyway (see
        Section~\ref{sec:userland-build}).
  \item \textbf{A rich standard library.} Packages like \code{os}, \code{io},
        \code{io/fs}, \code{text/tabwriter}, \code{archive/tar}, and
        \code{compress/gzip} provide most of what a coreutil needs out of the
        box, while \code{syscall} provides an escape hatch to raw ioctls when a
        utility must talk to the kernel directly.
\end{itemize}

\subsection{The Multicall (BusyBox-Style) Design}

A naive approach would compile each utility into its own binary. With Go that is
disastrous for disk footprint: every Go executable embeds the entire runtime and
whatever standard-library packages it imports, so fifteen separate binaries would
each carry several megabytes of duplicated runtime. The AULinux userland instead
follows the \emph{multicall} pattern popularised by BusyBox: \textbf{one} binary
contains \textbf{all} the commands, and the program decides which command to run
by inspecting the name it was invoked under.

\begin{infobox}[The single-binary payoff]
The whole userland---\code{ls}, \code{cat}, \code{cp}, \code{mv}, \code{mkdir},
\code{rm}, \code{chmod}, \code{echo}, \code{pwd}, \code{ps}, \code{ip},
\code{useradd}, \code{groupadd}, \code{aubuild}, \emph{and} the \code{aush}
shell---compiles into one statically linked binary named \code{auutils}. On the
root filesystem, \file{/bin/ls}, \file{/bin/cat}, \file{/bin/aush}, and the rest
are simply symlinks pointing at \code{auutils}. The Go runtime is paid for once.
\end{infobox}

The build script makes this concrete. It compiles the multicall binary and then
creates a ``link farm'' of symlinks, one per command, all pointing at the single
executable:

\begin{lstlisting}[language=bash, caption={The multicall build and symlink farm (scripts/build.sh)}]
CGO_ENABLED=0 go build \
    -ldflags="-s -w -buildid= -extldflags '-static'" \
    -trimpath -o "$BUILD_DIR/utils/auutils" utils/main.go

UTILS=("ls" "cat" "cp" "mv" "mkdir" "rm" "chmod" "echo" "pwd" \
       "ps" "ip" "useradd" "groupadd" "aubuild" "aush" "sh")
for tool in "${UTILS[@]}"; do
    ln -sf auutils "$BUILD_DIR/utils/$tool"
done
\end{lstlisting}

\section{The Multicall Dispatcher}

The heart of the userland is \file{utils/main.go}. It is the only
\code{package main} in the tree; every individual command and the shell are
ordinary library packages that expose a single exported entry point. The
dispatcher's job is purely to decide \emph{which} entry point to call.

\subsection{argv Routing}

When the kernel \code{exec}s one of the symlinks, \code{os.Args[0]} carries the
name the binary was invoked as---for example \code{/bin/ls}. The dispatcher
takes the basename of that path and switches on it. If the binary is instead
invoked under its own name, \code{auutils}, the first command-line argument is
treated as the command, exactly like \code{busybox ls -l}.

\begin{lstlisting}[language=Go, caption={argv-based command resolution (utils/main.go)}]
func main() {
    // Determine which tool to run based on argv[0]
    name := filepath.Base(os.Args[0])

    // If called as the multicall binary itself (e.g., "auutils"),
    // check if the first argument is a command.
    args := os.Args[1:]
    if name == "auutils" {
        if len(args) > 0 {
            name = args[0]
            args = args[1:]
        } else {
            // Default behavior: run shell if no args
            name = "aush"
        }
    }
    // ... switch (name) { ... }
}
\end{lstlisting}

Two design choices stand out. First, when \code{auutils} is run with no
arguments at all, it falls back to launching the shell---a friendly default that
makes the bare binary usable interactively. Second, the routing is a plain
\code{switch} statement, so there is no reflection, no init-time registry magic,
and the set of commands is statically knowable at compile time.

\begin{lstlisting}[language=Go, caption={The dispatch switch (utils/main.go)}]
switch name {
case "ls":
    ls.Run(args)
case "cat":
    cat.Run(args)
case "cp":
    cp.Run(args)
// ... mv, mkdir, rm, chmod, echo, pwd, ps, ip, useradd, groupadd, aubuild ...
case "aush", "sh": // "sh" is often a symlink to the shell
    shell.Run(args)
default:
    fmt.Fprintf(os.Stderr, "auutils: command not found: %s\n", name)
    os.Exit(1)
}
\end{lstlisting}

Note that both \code{aush} and the conventional name \code{sh} route to the same
shell, so scripts that hard-code \code{/bin/sh} keep working.

\subsection{The \code{Run(args)} Convention}

Every command package follows the same contract: it exposes a single exported
function \code{func Run(args []string)} that receives the argument slice
\emph{after} the command name has been stripped. The package itself is named
after the command (\code{package ls}, \code{package cp}, and so on)---not
\code{main}---which is precisely what makes them importable into the single
multicall binary.

\begin{notebox}[Verified: this is a true multicall binary, not separate binaries]
Despite each utility living in its own directory under \file{utils/cmd/<name>/main.go},
none of those files declare \code{package main}. They are library packages
(\code{package ls}, \code{package cp}, \dots) imported by \file{utils/main.go}.
The git history confirms a deliberate refactor from standalone binaries into this
multicall design ``to reduce disk footprint.'' The on-disk commands are symlinks
to \code{auutils}; there is exactly one executable.
\end{notebox}

Each \code{Run} is responsible for its own flag parsing (most use the standard
\code{flag} package with an isolated \code{flag.NewFlagSet}), its own usage text,
and---importantly---its own exit code via \code{os.Exit}. Because the utilities
call \code{os.Exit} directly on error, a non-zero status terminates the whole
process, which is exactly the semantics a shell expects from an external command.

\section{A Tour of \code{auutils}}

The multicall binary implements fourteen commands plus the shell. The table
below summarises each one and the behaviour that is actually present in the
source---not aspirational flags, but the ones wired up in code.

\begin{table}[h]
\centering
\small
\begin{tabular}{@{}lll@{}}
\toprule
\textbf{Command} & \textbf{Purpose} & \textbf{Notable behaviour implemented} \\
\midrule
\code{ls}       & List directory contents & \code{-l -a -A -h -t -S -r}, ANSI colour, TTY detection \\
\code{cat}      & Concatenate files       & \code{-n -b -E -T -s -v}, stdin via \code{-} \\
\code{cp}       & Copy files/dirs         & \code{-r/-R -f -i -v -p -n}, files, dirs, symlinks \\
\code{mv}       & Move/rename             & \code{-f -i -v -n}, atomic rename + cross-device fallback \\
\code{mkdir}    & Create directories      & \code{-p -m -v} \\
\code{rm}       & Remove files/dirs       & \code{-f -i -r/-R -v -d} \\
\code{chmod}    & Change mode bits        & \code{-R -v}, \emph{octal modes only} \\
\code{echo}     & Print text              & \code{-n -e -E}, backslash escapes incl. octal \\
\code{pwd}      & Print working directory & \code{-L -P} \\
\code{ps}       & List processes          & Parses \file{/proc/<pid>/stat}; PID/PPID/STATE/CMD \\
\code{ip}       & Interface up/down       & \code{ip link set <dev> [up|down]} via raw ioctl \\
\code{useradd}  & Add a user              & Appends to \file{/etc/passwd}, auto-allocates UID \\
\code{groupadd} & Add a group             & Appends to \file{/etc/group}, auto-allocates GID \\
\code{aubuild}  & Build a package         & Produces \code{.json} metadata + \code{.pkg} tar.gz \\
\bottomrule
\end{tabular}
\caption{Commands implemented in the \code{auutils} multicall binary.}
\end{table}

We now deep-dive into four representative commands that show the breadth of
techniques used: \code{ls} (filesystem metadata and formatting), \code{cp}
(recursive traversal and symlink handling), \code{ps} (parsing the
\file{/proc} pseudo-filesystem), and \code{ip} (raw syscalls).

\subsection{\code{ls}: Metadata, Sorting, and Colour}

\code{ls} is the most feature-rich utility. It supports long format
(\code{-l}), hidden-file display (\code{-a}/\code{-A}), human-readable sizes
(\code{-h}), and three sort orders---by name (default), modification time
(\code{-t}), or size (\code{-S})---each reversible with \code{-r}. Colour is
applied only when standard output is a terminal, detected by checking the
character-device bit on the file's mode:

\begin{lstlisting}[language=Go, caption={Terminal detection gates colour output (utils/cmd/ls/main.go)}]
func isTerminal(f *os.File) bool {
    fi, err := f.Stat()
    if err != nil {
        return false
    }
    return (fi.Mode() & os.ModeCharDevice) != 0
}
// ...
isTTY := isTerminal(os.Stdout)
useColor := isTTY && !*noColor
\end{lstlisting}

This is the correct, pipe-friendly behaviour: \code{ls -l | cat} emits no escape
codes because the destination is a pipe, not a character device. For the long
listing, ownership is resolved from the raw \code{Stat\_t} returned by
\code{info.Sys()}, mapping numeric UID/GID back to names and gracefully falling
back to the number when no name exists:

\begin{lstlisting}[language=Go, caption={Resolving owner and group from syscall.Stat\_t (utils/cmd/ls/main.go)}]
if stat, ok := info.Sys().(*syscall.Stat_t); ok {
    if u, err := user.LookupId(strconv.Itoa(int(stat.Uid))); err == nil {
        owner = u.Username
    } else {
        owner = strconv.Itoa(int(stat.Uid))
    }
    if g, err := user.LookupGroupId(strconv.Itoa(int(stat.Gid))); err == nil {
        group = g.Name
    } else {
        group = strconv.Itoa(int(stat.Gid))
    }
}
\end{lstlisting}

Colour selection is driven by file type and extension: directories are blue,
executables (any of the \code{0111} permission bits set) green, symlinks cyan,
archives red, and images magenta. Columns in long mode are aligned with
\code{text/tabwriter}, which removes the manual padding arithmetic that makes the
C original error-prone.

\subsection{\code{cp}: Recursion, Symlinks, and Preservation}

\code{cp} dispatches on the source type. The \code{copyPath} function is the
recursion pivot: directories require \code{-r} and recurse into \code{copyDir};
symlinks are recreated as symlinks (not dereferenced); everything else is a
regular file copy.

\begin{lstlisting}[language=Go, caption={Type-directed copy dispatch (utils/cmd/cp/main.go)}]
func copyPath(src, dest string) error {
    srcInfo, err := os.Lstat(src)
    if err != nil {
        return err
    }
    if srcInfo.IsDir() {
        if !*recursive {
            return fmt.Errorf("omitting directory '%s'", src)
        }
        return copyDir(src, dest, srcInfo)
    }
    if srcInfo.Mode()&os.ModeSymlink != 0 {
        return copySymlink(src, dest)
    }
    return copyFile(src, dest, srcInfo)
}
\end{lstlisting}

The use of \code{os.Lstat} (rather than \code{os.Stat}) is deliberate: it does
not follow the final symlink, which is what lets \code{copySymlink} recreate the
link verbatim with \code{os.Readlink}/\code{os.Symlink}. File contents are
streamed with \code{io.Copy}, so even very large files are copied with a bounded
buffer rather than read fully into memory. The \code{-p} flag preserves the
source mode; without it, the destination mode is masked to the permission bits
\code{0777}. The flags \code{-i} (interactive), \code{-n} (no-clobber), and
\code{-f} (force) all interact at the point where the destination already exists.

\subsection{\code{ps}: Reading the \texttt{/proc} Filesystem}

\code{ps} carries no kernel coupling beyond the standard Linux \file{/proc}
contract. It walks \file{/proc}, treats any directory whose name parses as an
integer as a PID, and reads that process's \file{stat} file. The subtle part is
parsing the \code{comm} field, which is wrapped in parentheses and may itself
contain spaces or parentheses; the code locates the \emph{last} \code{)} to
delimit it safely:

\begin{lstlisting}[language=Go, caption={Robustly parsing /proc/<pid>/stat (utils/cmd/ps/main.go)}]
content := string(data)
startComm := strings.Index(content, "(")
endComm := strings.LastIndex(content, ")")
if startComm == -1 || endComm == -1 || endComm < startComm {
    continue
}
comm := content[startComm+1 : endComm]
rest := content[endComm+2:] // Skip ") "
fields := strings.Fields(rest)
if len(fields) < 2 {
    continue
}
state := fields[0]
ppid, _ := strconv.Atoi(fields[1])
\end{lstlisting}

The scanning logic is factored into \code{scanProc(root string)} which takes the
\file{/proc} root as a parameter rather than hard-coding it---a small design
choice that makes the function unit-testable (Section~\ref{sec:userland-testing}).
Output is rendered through \code{text/tabwriter} as four aligned columns: PID,
PPID, STATE, CMD.

\begin{notebox}[\code{ps} is intentionally minimal]
There is no flag parsing in \code{ps}---it ignores its arguments entirely and
always prints the same four columns for every process. There is no \code{-e},
\code{-f}, \code{aux}, user-name resolution, or CPU/memory accounting. It is a
clean, correct snapshot of the process table and nothing more.
\end{notebox}

\subsection{\code{ip}: Talking to the Kernel via ioctl}

\code{ip} is the one utility that reaches past the standard library into raw
system calls. It implements exactly one operation---\code{ip link set <dev> up}
or \code{down}---by opening a datagram socket and issuing \code{SIOCGIFFLAGS} /
\code{SIOCSIFFLAGS} ioctls against an \code{ifreq} structure, the same ABI the
real \code{iproute2} uses.

\begin{lstlisting}[language=Go, caption={Bringing an interface up with a raw ioctl (utils/cmd/ip/main.go)}]
fd, err := syscall.Socket(syscall.AF_INET, syscall.SOCK_DGRAM, 0)
// ...
var ifr ifreq
copy(ifr.Name[:], ifName)

// Get current flags
syscall.Syscall(syscall.SYS_IOCTL, uintptr(fd),
    SIOCGIFFLAGS, uintptr(unsafe.Pointer(&ifr)))

if action == "up" {
    ifr.Flags |= (IFF_UP | IFF_RUNNING)
} else if action == "down" {
    ifr.Flags &^= (IFF_UP | IFF_RUNNING)
}
// Set new flags
syscall.Syscall(syscall.SYS_IOCTL, uintptr(fd),
    SIOCSIFFLAGS, uintptr(unsafe.Pointer(&ifr)))
\end{lstlisting}

The \code{ifreq} Go struct is laid out to match the kernel's \code{struct ifreq}
(a 16-byte name field, the flags, and padding), and \code{unsafe.Pointer} bridges
it into the syscall. With no action argument, the command instead just prints the
interface's current flags. This is a deliberately small slice of the real
\code{ip} command---enough to bring \code{lo} or \code{eth0} up at boot.

\begin{warnbox}[Privilege and scope]
\code{ip}, \code{useradd}, and \code{groupadd} all require root. \code{ip} needs
privileges to set interface flags; \code{useradd}/\code{groupadd} open
\file{/etc/passwd} and \file{/etc/group} for appending. None of them perform
locking on those files, so concurrent invocations could race. They are designed
for single-user administrative use during system setup, not high-concurrency
provisioning.
\end{warnbox}

\subsection{\code{aubuild}: The Package Builder}

\code{aubuild} is the userland's bridge to the Rust package manager covered in
Chapter~5. Given a source directory it walks the tree to compute the installed
size, emits a JSON manifest whose field tags mirror \code{pkg-manager/src/model.rs},
and produces a gzip-compressed tar archive with the \code{.pkg} extension using
\code{archive/tar} and \code{compress/gzip}. The archive paths are made relative
to the source root so the package is relocatable.

\section{\code{aush}: The Absolutely Unique Shell}

The interactive shell lives in \file{shell/shell.go} as \code{package shell},
exposing the same \code{Run(args)} entry point as the utilities so the dispatcher
can launch it uniformly. It supports three modes of operation: an interactive
REPL, single-command execution via \code{-c}, and non-interactive script file
execution. The shell is a self-contained implementation with builtins, pipes, I/O
redirection, environment-variable expansion, background jobs, command history,
and a coloured prompt.

\subsection{Shell State and Job Model}

The shell's state is a single \code{Shell} struct carrying the current user,
hostname, history buffer, last exit code, and the job table. Jobs are modelled
explicitly with a four-state lifecycle.

\begin{lstlisting}[language=Go, caption={The job and shell state model (shell/shell.go)}]
type JobState int
const (
    JobRunning JobState = iota
    JobStopped
    JobDone
    JobTerminated
)

type Job struct {
    ID       int
    PID      int
    Command  string
    State    JobState
    Process  *os.Process
    DoneChan chan struct{}
}

type Shell struct {
    user       *user.User
    hostname   string
    history    []string
    historyPos int
    exitCode   int
    running    bool
    loginShell bool
    jobs       []*Job
    nextJobID  int
}
\end{lstlisting}

Each background job carries a \code{DoneChan} that a goroutine closes when the
process exits---a clean, idiomatic Go way to signal completion that the
\code{fg} builtin waits on.

\subsection{Initialisation and Environment Bootstrap}

\code{NewShell} initialises the \code{Shell} struct, resolves the current user
via \code{user.Current()}, and reads the hostname. Critically, it also
\emph{bootstraps essential environment variables} that may be absent on a freshly
booted minimal system---before any child process is ever spawned:

\begin{lstlisting}[language=Go, caption={Environment variable bootstrap in NewShell() (shell/shell.go)}]
// Ensure essential environment variables are set
if os.Getenv("PATH") == "" {
    os.Setenv("PATH", "/bin:/sbin:/usr/bin:/usr/sbin")
}
if os.Getenv("HOME") == "" {
    os.Setenv("HOME", "/root")
}
if os.Getenv("USER") == "" {
    os.Setenv("USER", "root")
}
if os.Getenv("SHELL") == "" {
    os.Setenv("SHELL", "/bin/aush")
}
\end{lstlisting}

This four-variable bootstrap is the correct minimal set for a functioning POSIX
environment. \code{PATH} is the most critical: without it, \code{exec.LookPath}
finds nothing and every external command fails. Setting \code{HOME} to
\code{/root} ensures history persistence and \code{cd} without arguments work
even before login infrastructure is in place. \code{USER} and \code{SHELL} are
set so that scripts and programs that read them get sensible defaults rather than
empty strings. None of these assignments overwrite an existing value, so a
properly initialised login environment is never disturbed.

After the environment is guaranteed, \code{NewShell} calls \code{loadHistory} to
restore persisted history from \code{\textasciitilde/.aush\_history}.

\subsection{The \code{Run(args)} Entry Point and Execution Modes}

\code{Run} is the multicall entry point. It parses the argument slice and
branches into one of three execution paths:

\begin{itemize}
  \item \code{-c "command"} --- executes a single command string and exits
        immediately with that command's exit code.
  \item \code{scriptfile} --- a positional (non-flag) argument names a script
        file; \code{RunScript} is called and its return value passed to
        \code{os.Exit}.
  \item (no arguments) --- drops into the interactive REPL via \code{Start}.
\end{itemize}

\begin{lstlisting}[language=Go, caption={Argument parsing and mode dispatch in Run() (shell/shell.go)}]
func Run(args []string) {
    shell := NewShell()
    scriptFile := ""

    for i := 0; i < len(args); i++ {
        arg := args[i]
        switch arg {
        case "-l", "--login":
            shell.loginShell = true
        case "-c":
            if i+1 < len(args) {
                os.Exit(shell.ExecuteCommand(args[i+1]))
            }
            fmt.Fprintln(os.Stderr, "aush: -c requires an argument")
            os.Exit(1)
        case "--version":
            fmt.Printf("aush (AULinux Shell) %s\n", version)
            os.Exit(0)
        case "--help":
            printHelp()
            os.Exit(0)
        default:
            if !strings.HasPrefix(arg, "-") && scriptFile == "" {
                scriptFile = arg
            }
        }
    }

    if scriptFile != "" {
        os.Exit(shell.RunScript(scriptFile))
    }

    os.Exit(shell.Start())
}
\end{lstlisting}

The script-file path is detected by the \code{default} branch: any argument that
does not begin with \code{-} and is the first such argument is treated as a
filename. This matches the POSIX convention of \code{sh script.sh}.

\subsection{Non-Interactive Mode: \code{RunScript}}

\code{RunScript} implements batch execution of a shell script file. It opens the
file, scans it line by line with \code{bufio.Scanner}, skips blank lines and
comment lines (those starting with \code{\#}), and calls \code{s.ExecuteCommand}
for each remaining line. The return value is the exit code of the last command
executed, or \code{1} on any file or scan error.

\begin{lstlisting}[language=Go, caption={Script file execution via RunScript (shell/shell.go)}]
func (s *Shell) RunScript(filename string) int {
    file, err := os.Open(filename)
    if err != nil {
        fmt.Fprintf(os.Stderr, "aush: %v\n", err)
        return 1
    }
    defer file.Close()

    scanner := bufio.NewScanner(file)
    lastExit := 0
    for scanner.Scan() {
        line := strings.TrimSpace(scanner.Text())
        // Skip empty lines and comments
        if line == "" || strings.HasPrefix(line, "#") {
            continue
        }
        lastExit = s.ExecuteCommand(line)
    }
    if err := scanner.Err(); err != nil {
        fmt.Fprintf(os.Stderr, "aush: error reading script %s: %v\n", filename, err)
        return 1
    }
    return lastExit
}
\end{lstlisting}

A script can therefore contain any command the interactive shell understands:
external programs, builtins, pipes, and redirections. The shebang line
(\code{\#!/bin/aush}) is naturally ignored because it begins with \code{\#}.
What is \emph{not} supported in script mode is identical to what is missing from
the interactive shell: there are no \code{if}/\code{for}/\code{while} control
structures, no variable assignment (\code{VAR=value}), and no arithmetic
expansion. Scripts are sequential command lists---closer to a \code{.sh} batch
file than a full POSIX shell script.

\begin{notebox}[Script mode is a meaningful step forward]
Before this change, \code{aush} had no way to drive a script file: the only
non-interactive path was \code{-c}, which accepts exactly one command. With
\code{RunScript} in place, multi-step boot sequences, build recipes, and
configuration scripts can be written as plain text files and executed directly
with \code{aush script.sh} or \code{/bin/aush script.sh} from the shebang.
The limitation---no control flow---is real, but the set of things that
\emph{can} be scripted (bring up an interface, run the package manager, create
directories, set environment variables via \code{export}) covers the core
AULinux use-cases.
\end{notebox}

\subsection{The Read--Eval--Print Loop}

\code{Start} runs the interactive loop: print prompt, read a line, trim it,
record it in history, execute it, and repeat until \code{running} becomes false.

\begin{lstlisting}[language=Go, caption={The main REPL loop (shell/shell.go)}]
reader := bufio.NewReader(os.Stdin)
for s.running {
    s.printPrompt()
    line, err := reader.ReadString('\n')
    if err != nil {
        if err.Error() == "EOF" {
            fmt.Println()
            break
        }
        continue
    }
    line = strings.TrimSpace(line)
    if line == "" {
        continue
    }
    s.addHistory(line)
    s.exitCode = s.ExecuteCommand(line)
}
s.saveHistory()
\end{lstlisting}

\subsection{Tokenising and Parsing}

\code{ExecuteCommand} is the dispatcher for a single line. Its first decision is
whether the line contains a pipe; if so it hands off to \code{executePipeline}.
Otherwise it tokenises the line, detects a trailing \code{\&} for background
execution, applies redirections, expands variables, and finally runs either a
builtin or an external command.

Tokenising is done by \code{parseCommand}, a hand-written scanner---not a regex
split---that understands single and double quotes and backslash escapes:

\begin{lstlisting}[language=Go, caption={Quote- and escape-aware tokeniser (shell/shell.go)}]
for i := 0; i < len(line); i++ {
    ch := line[i]
    if inQuote {
        if ch == quoteChar {
            inQuote = false
        } else {
            current.WriteByte(ch)
        }
    } else {
        switch ch {
        case '"', '\'':
            inQuote = true
            quoteChar = ch
        case ' ', '\t':
            if current.Len() > 0 {
                args = append(args, current.String())
                current.Reset()
            }
        case '&':
            // split '&' off as its own token
        case '\\':
            if i+1 < len(line) {
                i++
                current.WriteByte(line[i])
            }
        default:
            current.WriteByte(ch)
        }
    }
}
\end{lstlisting}

\begin{notebox}[Quoting is simple but not full POSIX]
Quotes are stripped and whitespace inside them is preserved, which covers the
common cases. However there is no distinction between single- and double-quote
expansion semantics: variable expansion runs over all tokens \emph{after}
parsing, so \code{'\$HOME'} in single quotes will still be expanded. There is
also no glob expansion (\code{*}, \code{?}) anywhere in the shell.
\end{notebox}

\subsection{Variable Expansion}

After tokenising, \code{expandVariables} runs each argument through
\code{os.Expand}, with special handling for the standard shell pseudo-variables:

\begin{lstlisting}[language=Go, caption={Environment and special-variable expansion (shell/shell.go)}]
result[i] = os.Expand(arg, func(key string) string {
    switch key {
    case "?":
        return fmt.Sprintf("%d", s.exitCode)
    case "$":
        return fmt.Sprintf("%d", os.Getpid())
    case "HOME":
        if s.user != nil { return s.user.HomeDir }
    case "USER":
        if s.user != nil { return s.user.Username }
    case "SHELL":
        return os.Args[0]
    }
    return os.Getenv(key)
})
\end{lstlisting}

So \code{\$?} yields the last exit status, \code{\$\$} the shell's PID, and any
other \code{\$NAME} falls through to the process environment.

\subsection{Builtins versus External Commands}

Builtins are registered in a package-level map initialised in \code{init()}.
After expansion, \code{ExecuteCommand} checks whether the first token names a
builtin; if so it calls the handler directly, otherwise it execs an external
program.

\begin{lstlisting}[language=Go, caption={The builtin registry (shell/shell.go)}]
builtins = map[string]func(*Shell, []string) int{
    "cd":   builtinCd,      "exit":    builtinExit,
    "export": builtinExport, "unset":  builtinUnset,
    "pwd":  builtinPwd,     "echo":    builtinEcho,
    "history": builtinHistory, "help": builtinHelp,
    "version": builtinVersion, "type": builtinType,
    "jobs": builtinJobs,    "fg":      builtinFg,
    "bg":   builtinBg,
}
\end{lstlisting}

The builtins cover the cases that \emph{must} run inside the shell's own process
(because they mutate shell state): \code{cd} changes the process working
directory and understands \code{\~} expansion; \code{export}/\code{unset} mutate
the environment that child processes inherit; \code{exit} sets \code{running}
false and parses an optional exit code; \code{jobs}/\code{fg}/\code{bg} drive the
job table; and \code{type} reports whether a name is a builtin or resolves it on
\code{PATH} via \code{exec.LookPath}.

External commands are run through \code{os/exec}. The non-pipeline,
non-background path is the simplest---construct the command, wire up the
(possibly redirected) standard streams, and run it, translating Go's
\code{*exec.ExitError} back into a numeric exit code:

\begin{lstlisting}[language=Go, caption={Executing an external command (shell/shell.go)}]
cmd := exec.Command(args[0], args[1:]...)
cmd.Stdin = stdin
cmd.Stdout = stdout
cmd.Stderr = stderr
// ...
if err := cmd.Run(); err != nil {
    if exitErr, ok := err.(*exec.ExitError); ok {
        return exitErr.ExitCode()
    }
    fmt.Fprintf(os.Stderr, "aush: %s: %v\n", args[0], err)
    return 127
}
return 0
\end{lstlisting}

The exit code \code{127}---``command not found''---matches the convention every
other shell uses, so scripts checking \code{\$?} behave as expected.

\subsection{Redirection and Pipes}

I/O redirection is implemented in \code{parseRedirections}, which scans the token
list, consumes redirection operators, and opens the appropriate files. The
supported operators are \code{<} (stdin), \code{>} (truncating stdout),
\code{>>} (appending stdout), \code{2>} (stderr to a file), and the special
\code{2>\&1} (merge stderr into stdout). Opened files are closed via a
\code{defer} after the command completes.

Pipelines are genuinely implemented---not faked. \code{executePipeline} splits on
\code{|}, builds one \code{*exec.Cmd} per stage, and chains them by connecting
each stage's \code{StdoutPipe()} to the next stage's \code{Stdin}:

\begin{lstlisting}[language=Go, caption={Wiring a multi-stage pipeline (shell/shell.go)}]
// Connect pipes
for i := 0; i < len(cmds)-1; i++ {
    pipe, err := cmds[i].StdoutPipe()
    if err != nil {
        fmt.Fprintf(os.Stderr, "aush: pipe error: %v\n", err)
        return 1
    }
    cmds[i+1].Stdin = pipe
}
// First command reads stdin, last writes stdout
cmds[0].Stdin = os.Stdin
cmds[len(cmds)-1].Stdout = os.Stdout
cmds[len(cmds)-1].Stderr = os.Stderr

for _, cmd := range cmds {
    if err := cmd.Start(); err != nil { /* ... */ }
}
\end{lstlisting}

All stages are started, then waited on; the pipeline's exit status is taken from
the commands as they are reaped.

\begin{warnbox}[Pipelines route through external binaries only]
A subtle limitation: \code{executePipeline} builds every stage with
\code{exec.Command}, so \emph{builtins cannot participate in a pipeline}.
\code{echo hi | cat} works because \code{echo} also exists as an external
\code{auutils} command, but a pipeline stage is never matched against the
builtin map. The single-command path is the only one that consults builtins.
\end{warnbox}

\subsection{Background Jobs and Job Control}

A trailing \code{\&} marks a command for background execution. The shell starts
the process, records a \code{Job}, prints the \code{[id] pid} notification, and
spawns a goroutine that waits for the process and closes the job's
\code{DoneChan} on completion:

\begin{lstlisting}[language=Go, caption={Launching and tracking a background job (shell/shell.go)}]
job := &Job{
    ID:       s.nextJobID,
    PID:      cmd.Process.Pid,
    Command:  strings.Join(args, " "),
    State:    JobRunning,
    Process:  cmd.Process,
    DoneChan: make(chan struct{}),
}
s.jobs = append(s.jobs, job)
s.nextJobID++
fmt.Printf("[%d] %d\n", job.ID, job.PID)
go func() {
    cmd.Wait()
    job.State = JobDone
    close(job.DoneChan)
    fmt.Printf("[%d] Done %s\n", job.ID, job.Command)
}()
\end{lstlisting}

The \code{jobs} builtin lists active jobs (after pruning completed ones),
\code{bg} resumes a stopped job by sending \code{SIGCONT} via
\code{job.Process.Signal}, and \code{fg} blocks on the job's \code{DoneChan}.

\begin{notebox}[Job control is partial by design]
The shell tracks jobs and can continue them with \code{SIGCONT}, but it does
\emph{not} perform terminal process-group management. The \code{bringToForeground}
function notes in a comment that ``In a real shell, we would \code{tcsetpgrp}
here''---it simply waits on the done channel instead. There is likewise no
\code{SIGTSTP}/Ctrl-Z handler to move a foreground job into the \code{JobStopped}
state, so that state is defined but never reached in normal use.
\end{notebox}

\subsection{Signals, History, and the Prompt}

Signal handling is deliberately light. \code{setupSignals} installs a handler for
\code{SIGINT} and \code{SIGTERM}: on \code{SIGINT} (Ctrl-C) it prints a newline
and redraws the prompt rather than killing the shell, and on \code{SIGTERM} it
sets \code{running} false for a graceful exit.

History is persisted to \code{\textasciitilde/.aush\_history}, capped at 1000
entries, with consecutive duplicates suppressed.

The prompt is assembled in \code{printPrompt} with ANSI colours---green user,
blue host, cyan directory---and collapses the home directory to \code{\~}. The
prompt symbol is \code{\#} for root (UID 0) and \code{\$} otherwise, the
universal Unix convention. Username resolution prefers the live environment over
the potentially-nil \code{s.user} struct: \code{os.Getenv("USER")} is consulted
first; only if that is empty does the code fall back to \code{s.user.Username},
with the final fallback being the literal string \code{"user"}.

\begin{lstlisting}[language=Go, caption={Prompt username resolution (shell/shell.go)}]
username := os.Getenv("USER")
if username == "" {
    if s.user != nil {
        username = s.user.Username
    } else {
        username = "user"
    }
}
\end{lstlisting}

This ordering matters: \code{s.user} is populated by \code{user.Current()},
which can fail or return a nil pointer in a stripped-down environment. The
environment variable, bootstrapped unconditionally by \code{NewShell}, is always
available and is therefore the reliable first choice.

\section{Testing}
\label{sec:userland-testing}

The userland includes a focused unit test for the trickiest parsing logic:
\file{utils/cmd/ps/ps\_test.go}. Rather than depending on the real, ever-changing
\file{/proc}, the test exploits the fact that \code{scanProc} takes its root
directory as a parameter. It builds a \emph{fake} \file{/proc} inside a
\code{t.TempDir()}, populating it with synthetic \code{stat} files, then asserts
that \code{scanProc} returns exactly the expected PIDs, PPIDs, command names, and
states.

\begin{lstlisting}[language=Go, caption={Testing /proc parsing against a synthetic tree (utils/cmd/ps/ps\_test.go)}]
tmpDir := t.TempDir()
// ... write tmpDir/<pid>/stat files like "100 (bash) S 1 ..." ...
// Create a non-pid directory (should be ignored)
os.Mkdir(filepath.Join(tmpDir, "sys"), 0755)

procs, err := scanProc(tmpDir)
if err != nil {
    t.Fatalf("scanProc failed: %v", err)
}
if len(procs) != len(pids) {
    t.Errorf("expected %d processes, got %d", len(pids), len(procs))
}
\end{lstlisting}

The test also seeds a non-numeric directory (\code{sys}) to confirm it is
ignored, validating the integer-name filter. This dependency-injection style---a
function taking its filesystem root as an argument---is the key enabler, and it is
a good template for testing the other utilities.

\begin{notebox}[Test coverage is currently narrow]
\code{ps} is the only command with a test file at present. The remaining
utilities and the shell are exercised manually (the git history references a
``comprehensive test script''), but they have no Go-level unit tests checked in.
This is an obvious area for future hardening.
\end{notebox}

\section{Build and Linking Notes}
\label{sec:userland-build}

The userland is one Go module (\file{go.mod} declares \code{module aulinux}),
which is why \file{utils/main.go} can import both \code{aulinux/shell} and every
\code{aulinux/utils/cmd/*} package into a single binary. The build flags chosen
in \file{scripts/build.sh} are all about producing the smallest possible static
artifact:

\begin{itemize}
  \item \code{CGO\_ENABLED=0} together with \code{-extldflags '-static'} yields a
        fully static binary with no libc dependency---it runs in an empty rootfs.
  \item \code{-ldflags="-s -w"} strips the symbol table and DWARF debug
        information, shaving a significant fraction off the binary size.
  \item \code{-trimpath} and \code{-buildid=} remove build-machine paths and the
        build ID for reproducible, leaner output.
  \item An optional \code{upx --best --lzma} pass compresses the executable
        further when UPX is available; the build degrades gracefully with a
        warning if it is not.
\end{itemize}

Runtime memory is tuned not in the userland itself but by \code{au-init}, which
sets \code{GOGC=50} and \code{GOMEMLIMIT=512MiB} in the environment before
launching the shell (see \file{init/src/init.c}). The lower \code{GOGC} makes the
collector run more aggressively, trading a little CPU for a smaller resident set
on memory-constrained systems---a sensible default for a userland of many
short-lived processes.

\section{Design Rationale and Limitations}

The Go userland is a pragmatic, honest piece of engineering: it implements the
commands a minimal system actually needs to boot, log in, manage files, inspect
processes, configure a network interface, and build packages---and it does so in
one static binary with real memory safety. The deep-dives above show genuine
attention to correctness: TTY-gated colour, \code{Lstat}-based symlink handling,
robust \file{/proc} parsing, and real pipelines.

It is equally important to be clear about the edges, all of which are visible in
the source:

\begin{itemize}
  \item \textbf{\code{chmod} accepts octal modes only.} A \code{// TODO} in the
        source notes that symbolic modes (\code{u+x}, \code{go-w}) are
        unimplemented; non-octal input is rejected with an error.
  \item \textbf{No globbing.} Neither the shell nor the utilities expand
        \code{*}/\code{?}; wildcards are passed through literally.
  \item \textbf{Builtins are excluded from pipelines.} Only the single-command
        path consults the builtin map; pipeline stages are always external.
  \item \textbf{Job control lacks terminal process-group handling.} \code{SIGCONT}
        works, but \code{tcsetpgrp} and Ctrl-Z stop/resume are not wired up;
        \code{JobStopped} and \code{JobTerminated} are defined but unreached.
  \item \textbf{No file locking in the account tools.} \code{useradd} and
        \code{groupadd} append to \file{/etc/passwd} and \file{/etc/group}
        without locking.
  \item \textbf{\code{mv}'s cross-device limits.} \code{mv} falls back to
        copy-and-delete across devices but explicitly refuses cross-device
        \emph{directory} moves.
  \item \textbf{Script mode lacks control flow.} \code{RunScript} executes
        commands sequentially and understands comments and blank lines, but there
        are no \code{if}/\code{for}/\code{while} constructs, no variable
        assignment (\code{VAR=value}), and no arithmetic expansion. Scripts are
        sequential command lists only.
  \item \textbf{Thin test coverage.} Only \code{ps} has a unit test today.
\end{itemize}

\begin{infobox}[The bigger picture]
None of these limitations block the system's core mission. AULinux's userland is
deliberately a \emph{focused} reimplementation, not a drop-in GNU coreutils
clone. By accepting a smaller, well-understood feature set, it stays small,
static, memory-safe, and entirely contained in a single auditable binary---which
is exactly the property a polyglot, from-scratch operating system wants from the
layer its users touch most.
\end{infobox}
